\documentclass[11pt,a4paper]{article}

\usepackage[margin=2.5cm]{geometry}
\usepackage{amsmath,amssymb,amsthm}
\usepackage{enumitem}
\usepackage{booktabs}
\usepackage{hyperref}
\usepackage{microtype}

\newtheorem{theorem}{Theorem}
\newtheorem{conjecture}[theorem]{Conjecture}
\newtheorem{proposition}[theorem]{Proposition}
\newtheorem{lemma}[theorem]{Lemma}
\newtheorem{corollary}[theorem]{Corollary}
\newtheorem{definition}[theorem]{Definition}
\newtheorem{example}[theorem]{Example}
\newtheorem{remark}[theorem]{Remark}

\newcommand{\Z}{\mathbb{Z}}
\newcommand{\N}{\mathbb{N}}
\newcommand{\Q}{\mathbb{Q}}
\newcommand{\F}{\mathbb{F}}
\newcommand{\Leg}[2]{\left(\frac{#1}{#2}\right)}

\title{Bunyakovsky's Conjecture:\\Prime Values of Irreducible Polynomials}
\author{Robin Langer}
\date{August 2026}

\begin{document}
\maketitle

\begin{abstract}
We give an expository account of Bunyakovsky's conjecture (1857), which asserts that every irreducible polynomial over $\Z$ with positive leading coefficient and no fixed prime divisor represents infinitely many primes. We trace the conjecture's position in the hierarchy of open problems in analytic number theory---from Dirichlet's theorem on primes in arithmetic progressions, through the Bateman--Horn conjecture's quantitative predictions, to the special case $f(x) = 2x^2 - 1$ arising from Chebyshev polynomials. We present computational evidence, explain the algebraic structure of the ``sieving obstructions,'' and describe the connection to Lucas sequences and primality testing.
\end{abstract}

\section{Introduction}

Which polynomials produce infinitely many prime values? The simplest case is well understood: Euclid proved there are infinitely many primes (the polynomial $f(x) = x$, in spirit), and Dirichlet's theorem (1837) establishes that $f(x) = ax + b$ represents infinitely many primes whenever $\gcd(a,b) = 1$. But for degree two and above, the question remains wide open.

In 1857, Viktor Yakovlevich Bunyakovsky proposed the natural generalisation \cite{bunyakovsky1857}.

\begin{conjecture}[Bunyakovsky, 1857]\label{conj:bun}
Let $f(x) \in \Z[x]$ be irreducible over $\Q$, with positive leading coefficient and $\deg f \geq 1$. Suppose $f$ has no \emph{fixed prime divisor}: there is no prime $p$ such that $p \mid f(n)$ for every $n \in \Z$. Then $f(n)$ is prime for infinitely many $n \in \N$.
\end{conjecture}

The condition on fixed prime divisors is necessary: the polynomial $f(x) = x^2 + x + 2$ is irreducible over $\Q$, yet $f(n)$ is always even, so it can be prime only at $n = 0$ (giving $f(0) = 2$). Without excluding such obstructions, the conjecture would be trivially false.

As of 2026, Conjecture~\ref{conj:bun} is unproven for \emph{every single polynomial of degree $\geq 2$}. Not a single quadratic, cubic, or higher-degree polynomial has been shown to produce infinitely many primes. The most famous special case---whether $n^2 + 1$ is prime infinitely often---remains open despite enormous effort.

\section{The Hierarchy of Conjectures}

Bunyakovsky's conjecture sits in a nested hierarchy, each level stronger than the last.

\begin{enumerate}[label=(\roman*)]
    \item \textbf{Dirichlet's theorem} (1837, proven): If $\gcd(a,b) = 1$, then $an + b$ is prime for infinitely many $n$.
    \item \textbf{Bunyakovsky's conjecture} (1857, open): Any irreducible $f \in \Z[x]$ without a fixed prime divisor represents infinitely many primes.
    \item \textbf{Schinzel's Hypothesis H} (1958, open): If $f_1, \ldots, f_k \in \Z[x]$ are irreducible with positive leading coefficients and no prime $p$ divides $\prod f_i(n)$ for all $n$, then $f_1(n), \ldots, f_k(n)$ are \emph{simultaneously} prime for infinitely many $n$ \cite{schinzel1958}.
    \item \textbf{Bateman--Horn conjecture} (1962, open): Under the same hypotheses as (iii), the number of $n \leq N$ for which all $f_i(n)$ are prime is asymptotically
    \begin{equation}\label{eq:bh}
        \pi_{f_1,\ldots,f_k}(N) \sim \frac{C_{f_1,\ldots,f_k}}{\prod_{i=1}^k \deg f_i} \cdot \frac{N}{(\ln N)^k},
    \end{equation}
    where $C_{f_1,\ldots,f_k}$ is an explicit product over primes \cite{bateman1962}.
\end{enumerate}

Each conjecture implies the ones above it. The Bateman--Horn conjecture is the strongest: it predicts not merely existence but an asymptotic count. When its predictions can be tested numerically, the agreement is striking.

\section{The Bateman--Horn Constant}

For a single polynomial $f$ of degree $d$, the Bateman--Horn prediction \eqref{eq:bh} becomes
\[
    \pi_f(N) \sim \frac{C_f}{d} \cdot \frac{N}{\ln N},
\]
where the constant is an Euler product:
\begin{equation}\label{eq:constant}
    C_f = \prod_{p \text{ prime}} \frac{1 - \omega_f(p)/p}{1 - 1/p},
\end{equation}
and $\omega_f(p)$ denotes the number of solutions to $f(x) \equiv 0 \pmod{p}$.

The intuition is probabilistic. If $n$ is a ``random'' integer near $N$, the prime number theorem says it is prime with probability $\sim 1/\ln N$. If instead we consider $f(n) \sim c \cdot n^d$ for large $n$, then $f(n)$ is prime with probability $\sim 1/(d \ln n)$. The Euler product $C_f$ corrects for local biases: if $f$ has more roots mod $p$ than a ``generic'' degree-$d$ polynomial, then $f$ is divisible by $p$ more often, reducing the expected prime count. The converse holds when $f$ has fewer roots mod $p$.

\begin{example}
For $f(x) = x^2 + 1$: the constant $C_f \approx 1.3728$ (the Landau--Ramanujan constant, up to normalization). The polynomial has $\omega_f(p) = 0$ when $p \equiv 3 \pmod{4}$ and $\omega_f(p) = 2$ when $p \equiv 1 \pmod{4}$, by quadratic reciprocity.
\end{example}

\section{Case Study: $f(x) = 2x^2 - 1$ and Chebyshev Polynomials}

The polynomial $f(x) = 2x^2 - 1$ is the Chebyshev polynomial of the first kind $T_2(x)$. It provides a particularly clean case study for Bunyakovsky's conjecture.

\subsection{Checking the hypotheses}

\begin{itemize}
    \item \emph{Irreducible over $\Q$}: The discriminant is $0^2 - 4(2)(-1) = 8$, which is not a perfect square, so $f$ is irreducible.
    \item \emph{Positive leading coefficient}: $2 > 0$.
    \item \emph{No fixed prime divisor}: $f(1) = 1$, $f(2) = 7$, $f(3) = 17$, so $\gcd(f(1), f(2), f(3)) = 1$.
\end{itemize}

Bunyakovsky's conjecture therefore predicts that $2n^2 - 1$ is prime for infinitely many $n$.

\subsection{The Bateman--Horn constant for $2x^2 - 1$}

We need $\omega_f(p)$, the number of solutions to $2x^2 \equiv 1 \pmod{p}$, i.e., the number of square roots of $2^{-1} \pmod{p}$.

\begin{proposition}\label{prop:roots}
Let $p$ be an odd prime. Then
\[
    \omega_f(p) = 1 + \Leg{2^{-1}}{p} = 1 + \Leg{(p+1)/2}{p},
\]
where $\Leg{\cdot}{p}$ is the Legendre symbol. In particular:
\[
    \omega_f(p) = \begin{cases}
        2 & \text{if } p \equiv \pm 1 \pmod{8}, \\
        0 & \text{if } p \equiv \pm 3 \pmod{8}.
    \end{cases}
\]
\end{proposition}

\begin{proof}
The equation $2x^2 \equiv 1 \pmod{p}$ has solutions iff $2^{-1}$ is a quadratic residue mod $p$. Since $\Leg{2^{-1}}{p} = \Leg{2}{p}^{-1} = \Leg{2}{p}$ (as the Legendre symbol takes values $\pm 1$), and the second supplement to quadratic reciprocity gives $\Leg{2}{p} = (-1)^{(p^2-1)/8}$, which equals $+1$ iff $p \equiv \pm 1 \pmod 8$.
\end{proof}

\begin{corollary}
The only primes that can divide $2n^2 - 1$ are $2$ (trivially: $2n^2 - 1$ is always odd) and primes $q \equiv \pm 1 \pmod{8}$. In particular, $2n^2 - 1$ is never divisible by $3$, $5$, $11$, $13$, $19$, $29$, or $37$.
\end{corollary}

This ``sieving'' by quadratic reciprocity explains why $2n^2 - 1$ produces primes at an unusually high rate: it is automatically coprime to a large set of small primes. The Bateman--Horn constant reflects this advantage. Computing the product \eqref{eq:constant} over primes $p < 500$:

\[
    C_f \approx 3.666.
\]

This is roughly 2.7 times larger than $C_{x^2+1} \approx 1.37$, predicting that $2n^2 - 1$ produces primes at nearly three times the rate of $n^2 + 1$.

\subsection{Computational evidence}

Table~\ref{tab:bh} compares the Bateman--Horn prediction to exact counts.

\begin{table}[h]
\centering
\begin{tabular}{rrrrr}
\toprule
$N$ & Predicted & Actual & Ratio \\
\midrule
100 & 39.8 & 45 & 1.131 \\
500 & 147.5 & 174 & 1.180 \\
1\,000 & 265.4 & 303 & 1.142 \\
5\,000 & 1\,076.1 & 1\,193 & 1.109 \\
10\,000 & 1\,990.1 & 2\,202 & 1.106 \\
50\,000 & 8\,470.5 & 9\,175 & 1.083 \\
\bottomrule
\end{tabular}
\caption{Count of $n \leq N$ with $2n^2 - 1$ prime: Bateman--Horn prediction vs.\ actual. The ratio converges toward 1, consistent with the conjecture.}
\label{tab:bh}
\end{table}

The ratio trends steadily toward 1 as $N$ grows, exactly as the conjecture predicts. The systematic overcount at small $N$ is typical: the Euler product converges slowly, and the truncation to primes $p < 500$ slightly underestimates $C_f$.

\subsection{Restricting to prime arguments}

When we evaluate $T_2(p) = 2p^2 - 1$ only at prime values of $p$, we ask a stronger question: is $2p^2 - 1$ prime for infinitely many primes $p$? Heuristically, there is no algebraic obstruction coupling the primality of $p$ to that of $2p^2 - 1$, so we expect the events to be approximately independent. Computationally, approximately 20\% of primes $p < 10\,000$ yield a prime $T_2(p)$.

\section{Connection to Lucas Sequences and Primality Testing}

The appearance of the Chebyshev polynomial $T_2$ is not a coincidence. There is a deep and exact identity connecting Chebyshev polynomials to Lucas sequences, and this identity underpins the most important primality tests in computational number theory.

\subsection{The Chebyshev--Lucas identity}

\begin{theorem}\label{thm:lucas}
Let $V_n(P, Q)$ denote the Lucas $V$-sequence defined by $V_0 = 2$, $V_1 = P$, $V_n = PV_{n-1} - QV_{n-2}$. Then for all $n \geq 0$ and all $P$:
\[
    V_n(P, 1) = 2 \, T_n\!\left(\frac{P}{2}\right),
\]
where $T_n$ is the $n$-th Chebyshev polynomial of the first kind.
\end{theorem}

\begin{proof}
Let $\alpha, \beta$ be the roots of $t^2 - Pt + 1 = 0$. Then $\alpha\beta = 1$, so $\beta = \alpha^{-1}$, and $\alpha + \alpha^{-1} = P$. Setting $\cos\theta = P/2$ (so $\alpha = e^{i\theta}$ when $|P| \leq 2$, or extending by analytic continuation):
\[
    V_n = \alpha^n + \alpha^{-n} = 2\cos(n\theta) = 2\,T_n(\cos\theta) = 2\,T_n(P/2). \qedhere
\]
\end{proof}

More generally, the \emph{Dickson polynomial} $D_n(x, a)$ satisfies $D_n(\alpha + a/\alpha, a) = \alpha^n + (a/\alpha)^n$, giving $V_n(P, Q) = D_n(P, Q)$ for all $Q$. Dickson polynomials are the natural generalisation of Chebyshev polynomials, and they inherit the crucial composition law:
\[
    D_m(D_n(x, a), a^n) = D_{mn}(x, a).
\]

\subsection{Lucas--Lehmer as iterated Chebyshev evaluation}

The Lucas--Lehmer primality test for Mersenne numbers $M_p = 2^p - 1$ proceeds as follows:
\begin{align*}
    s_0 &= 4, \\
    s_{i+1} &= s_i^2 - 2, \\
    M_p \text{ is prime} &\iff s_{p-2} \equiv 0 \pmod{M_p}.
\end{align*}

Each step of the recurrence is an application of $T_2$:
\[
    s_{i+1} = s_i^2 - 2 = 2T_2(s_i/2) = V_2(s_i, 1).
\]

By the composition law, after $k$ steps:
\[
    s_k = V_{2^k}(4, 1) = 2\,T_{2^k}(2).
\]

The Lucas--Lehmer test therefore asks a single question: \emph{does $M_p$ divide $T_{2^{p-2}}(2)$?} Primality testing, at its core, is iterated Chebyshev polynomial evaluation modulo a candidate.

\subsection{Why Chebyshev and not Hermite?}

The composition law $T_m \circ T_n = T_{mn}$ is unique to the Chebyshev/Dickson family among classical orthogonal polynomials. It arises because Chebyshev polynomials encode the power map $\alpha \mapsto \alpha^n$ on the unit circle, and composition of power maps corresponds to multiplication of exponents.

Hermite polynomials satisfy $H_{n+1}(x) = 2xH_n(x) - 2nH_{n-1}(x)$, where the coefficient $2n$ depends on the index. This breaks any hope of a composition law: the recurrence is not ``autonomous.'' Consequently, Hermite (and Legendre, Laguerre, etc.) do not encode the multiplicative structure of $\Z$, and their values at primes reflect only the divisibility properties of their coefficients---structured, but not deep.

\section{The Algebraic Structure of Obstructions}

Returning to Bunyakovsky's conjecture, we can ask: for a given irreducible $f$, what determines whether $f$ has a fixed prime divisor?

\begin{definition}
The \emph{fixed divisor} of $f \in \Z[x]$ is $\gcd\{f(n) : n \in \Z\}$.
\end{definition}

\begin{proposition}
For $f$ of degree $d$ with leading coefficient $a_d$:
\[
    \gcd\{f(n) : n \in \Z\} = \gcd\{f(0), f(1), \ldots, f(d)\}.
\]
In particular, the fixed divisor is computable from finitely many values.
\end{proposition}

When the fixed divisor is 1, there is no ``trivial'' reason for $f$ to avoid primes, and Bunyakovsky's conjecture asserts that it does not. The conjecture can be restated: \emph{the only obstructions to a polynomial representing primes are the local ones} (divisibility by a fixed prime). There are no ``global conspiracies.''

This philosophy---local obstructions are the only obstructions---pervades modern number theory. It appears in the Hasse--Minkowski theorem for quadratic forms, in the Brauer--Manin obstruction for rational points, and in the Hardy--Littlewood circle method. Bunyakovsky's conjecture is, in a sense, the assertion that this philosophy extends to prime values of polynomials.

\section{What Would a Proof Look Like?}

The fundamental difficulty is that we lack tools to control the prime factorisation of polynomial values. The strongest known result in this direction is:

\begin{theorem}[Iwaniec, 1978 \cite{iwaniec1978}]
There are infinitely many $n$ such that $n^2 + 1$ has at most two prime factors.
\end{theorem}

This is a sieve-theoretic result: Iwaniec shows that $n^2 + 1$ is ``almost prime'' infinitely often, but the sieve cannot distinguish between a prime and a product of two primes (the \emph{parity problem} in sieve theory).

A proof of Bunyakovsky's conjecture, even for a single degree-2 polynomial, would likely require fundamentally new ideas. Possible approaches include:
\begin{itemize}
    \item Resolving the parity barrier in sieve methods (a major open problem in its own right).
    \item Establishing sufficiently strong equidistribution results for primes in polynomial sequences, perhaps via automorphic forms or algebraic geometry.
    \item A conceptually different approach entirely---perhaps through the algebraic structure of function fields, where an analogue of Bunyakovsky's conjecture is known to hold.
\end{itemize}

\begin{remark}
Over function fields $\F_q[t]$, the analogue of Bunyakovsky's conjecture is a theorem: an irreducible polynomial $f(x) \in \F_q[t][x]$ represents infinitely many irreducible polynomials in $\F_q[t]$. The proof uses the Riemann hypothesis for curves over finite fields (Weil, 1948). The absence of a comparable tool over $\Z$ is the heart of the difficulty.
\end{remark}

\section{Concluding Remarks}

Bunyakovsky's conjecture, now 169 years old, remains as tantalising as ever. It occupies a peculiar position in number theory: universally believed, supported by overwhelming computational evidence, connected to deep structures in algebra and analysis, yet completely resistant to proof.

The case $f(x) = 2x^2 - 1$---the Chebyshev polynomial $T_2$---is especially appealing because it links the conjecture to primality testing via Lucas sequences. The same polynomial family that \emph{detects} primes (through Lucas--Lehmer and Baillie--PSW) also \emph{generates} them at a rate that we can predict but not prove. Whether this duality between detection and generation can be exploited toward a proof is an open and intriguing question.

\begin{thebibliography}{9}
\bibitem{bunyakovsky1857}
V.~Bunyakovsky, \emph{Sur les diviseurs num\'eriques invariables des fonctions rationnelles enti\`eres}, M\'em.\ Acad.\ Imp.\ Sci.\ St.-P\'etersbourg, \textbf{6} (1857), no.~7, 305--329.

\bibitem{bateman1962}
P.~T.~Bateman and R.~A.~Horn, \emph{A heuristic asymptotic formula concerning the distribution of prime numbers}, Math.\ Comp., \textbf{16} (1962), 363--367.

\bibitem{schinzel1958}
A.~Schinzel and W.~Sierpi\'nski, \emph{Sur certaines hypoth\`eses concernant les nombres premiers}, Acta Arith., \textbf{4} (1958), 185--208; erratum \textbf{5} (1959), 259.

\bibitem{iwaniec1978}
H.~Iwaniec, \emph{Almost-primes represented by quadratic polynomials}, Invent.\ Math., \textbf{47} (1978), 171--188.
\end{thebibliography}

\end{document}
